\documentclass[10pt,a4paper]{amsart}
\usepackage[margin=19mm]{geometry}
\usepackage{amsmath,amssymb,mathtools}
\usepackage[scr=rsfs,cal=euler]{mathalfa}
\usepackage{xfrac}
\usepackage[unicode,hidelinks,bookmarks=false]{hyperref}
\newcommand{\ie}{i.e.\ }
\newcommand{\cf}{cf.\ }
\newcommand{\resp}{respectively\ }
\newcommand{\Subsection}{\subsection}
\providecommand{\nobreakdash}{\nobreak\hbox{-}}
\setlength{\emergencystretch}{2em}
\multlinegap0pt
\usepackage{bm}
\usepackage{bbm}
\usepackage{dsfont}
\usepackage{xspace}
\usepackage{cancel}
\usepackage{mathtools}
\usepackage{amsthm}
\makeatletter
\@ifundefined{lemma}{\newtheorem{lemma}{Lemma}}{}
\@ifundefined{theorem}{\newtheorem{theorem}{Theorem}}{}
\makeatother
\def\m{\mathbf{m}}
\def\pt{\mathbf{p}}
\newcommand{\p}{{\mathbf{P}}}
\newcommand{\rr}{\mathbb{R}}
\title[Trees with extremal Laplacian eigenvalue multiplicity]{Trees with extremal Laplacian eigenvalue multiplicity}
\author{Vinayak Gupta, Gargi Lather, R. Balaji}
\date{Source: arXiv:2507.15472v1 (CC BY 4.0)}
\begin{document}
\maketitle

\noindent\emph{Source and licence.} Trees with extremal Laplacian eigenvalue multiplicity. Version arXiv:2507.15472v1 (CC BY 4.0), distributed under the Creative Commons Attribution 4.0 International licence, as recorded in the arXiv licence metadata for this exact version and shown on the abstract page https://arxiv.org/abs/2507.15472v1. Full rights notice: licenses/arxiv-2507.15472-CC-BY-4.0.txt.\par\smallskip

\noindent\textit{Editorial scope.} Counted unit: the Lemma whose source label is \texttt{lem6base} in the article (source0.tex lines 546--564, source label \texttt{lem6base}), delivered together with the article's own complete original proof of it (source0.tex lines 566--744, one \texttt{proof} environment of 179 lines). No mathematics is written, repaired or re-derived: statement and proof are the article's own text line for line. Required context retained uncounted: th1 (lines 303--315); lem1 (lines 321--329); v1 (lines 443--448); equiv (lines 513--519). Every definition, display and label of those lines is retained unchanged. Omitted from this package and not redistributed: 1--302, 316--320, 330--442, 449--512, 520--545, 565--565, 745--1308. Nothing inside a retained excerpt is omitted or altered: every retained body is delivered exactly as the article's lines are written, so no omission receipt of any kind accompanies this package. Citations: the retained excerpts carry no citation command at all, so no bibliography entry of the article is transcribed and no reference is redistributed. Reference adaptation: none was needed and none was made; every reference inside the delivered unit resolves inside it, so no unresolved reference or citation is produced. Printed slips kept literal: no printed word, symbol, hypothesis or number of the counted statement or of its proof was altered. Layout and class: the article's own class is \texttt{elsarticle} with packages including amssymb, amsmath, hyperref, amsthm; the wrapper is the lane's pinned \texttt{amsart} editorial wrapper at 10pt on A4, which restates the theorem-like environments the retained text needs and adds the article's own macro definitions that the retained text needs (\texttt{\textbackslash m}, \texttt{\textbackslash p}, \texttt{\textbackslash pt}, \texttt{\textbackslash rr}), with extra wrapper packages (bm, bbm, dsfont, xspace, cancel, mathtools). The delivered numbering is the unit's own. Assets: no retained span contains a figure environment, a graphics-inclusion command, a PDF-page inclusion, a table or a TikZ picture; no image file is copied into this package, the package declares no asset dependency and no image file is redistributed with it. Licence: version arXiv:2507.15472v1 of this article is distributed under the Creative Commons Attribution 4.0 International licence, as recorded in the arXiv licence metadata for this exact version and shown on the abstract page https://arxiv.org/abs/2507.15472v1. A full rights notice is delivered at licenses/arxiv-2507.15472-CC-BY-4.0.txt. Characters: every retained excerpt is pure ASCII, so the delivered engine, XeLaTeX with its default Latin Modern text font, reports no missing character.
\begin{theorem}\label{th1}
    Let $\p$ be the path $1 \sim 2 \sim \dotsc \sim n-1 \sim n$.
    Then,
    \(L_\p x^j=\lambda_j x^j,\)
    where 
    \begin{equation} \label{pn}
    \lambda_j:=2\left(1-\cos \frac{\pi j}{n}\right) ~~~j=0:n-1;
    \end{equation}
    and the vector $x^j=(x^j_{1},\dotsc,x^j_{n})'$ is given by
    \begin{equation} \label{xj}
    x^j_\nu:=\cos \left(\frac{\pi j}{n} \left (\nu-\frac{1}{2}\right ) \right) ~~~\nu=1:n.
    \end{equation}
\end{theorem}

\begin{lemma}\label{lem1}
Let $\p$ and $x_{\nu}^j$ be as in Theorem \ref{th1} and
$l \in \{0,\dotsc,n-1\}$ be such that
$\frac{l}{n}=\frac{2b+1}{2q+1}$,
where $b$ and $q$ are some integers.
If $u$ and $v$ are any two vertices of $\p$ such that
    $d(v,u) \equiv 0\,{\rm{mod}}\,(2q+1)$,
then $x^l_v=0$ if and only if $x^l_u=0$. 
\end{lemma}

\begin{lemma}\label{v1}
    Let $T_1$ and $T_2$ be trees with vertices $1,\dotsc,n$ and 
    $n,\dotsc,n+m$, respectively. If
  $y=(y_1,\dotsc,y_{n-1},0)'$ is a non-zero vector in $\rr^{n}$ such that
  $L_{T_1}y=\lambda y$, then $L_{T_{1} \circ T_2}(y_1,\dotsc,y_{n-1},0,\dotsc,0)'=\lambda (y_1,\dotsc,y_{n-1},0,\dotsc,0)'$.
\end{lemma}

\begin{lemma}\label{equiv}
  Let $T$ be a tree with at least one major vertex and $q$ be an integer. Then, the following are equivalent.
  \begin{enumerate}
      \item[\rm{(i)}] $d(u,w) \equiv 2q\,{\rm{mod}}\,(2q+1)$ for all distinct $u,w \in \pt(T)$.
      \item[\rm{(ii)}] $d(u,\theta) \equiv q\, {\rm{mod}}\, (2q+1)$ for all $u \in \pt(T)$ and $\theta \in \m(T)$.
  \end{enumerate}
  \end{lemma}

\begin{lemma}\label{lem6base}
Let $T$ be a tree with exactly one major vertex $v$
and $p$ pendant vertices.
Assume that there exists an integer $q$ 
such that
\[d(u_i,u_j)\equiv 2q\,{\rm{mod}}\,(2q+1)~~\mbox{for all distinct}~~u_i,u_j \in \pt(T).
\]
If $b$ is any integer such that $b \in [0,q)$, then
\[\lambda:=2\left(1- \cos \left(\frac{2b+1}{2q+1}\pi\right) \right)\]
is an eigenvalue of $L_T$. Furthermore,
there exist linearly independent vectors $x^1,\dots,x^{p-1}$
such that 
\begin{enumerate}
\item[\rm (a)] $L_Tx^j=\lambda x^j$ for all $j=1:p-1$.
\item[{\rm (b)}] %If $v$ is the major vertex, then
$x^j_{v}=0$ for all $j=1:p-1$. If $s \in V(T)$ and
$d(s,v) \equiv 0\,{\rm{mod}}\,(2q+1)$, then $x^j_{s}=0$ for all $j=1:p-1$. 
\end{enumerate}
\end{lemma}

\begin{proof}
 Let $n:=V(T)$. We prove the result by induction on $p$. Let $p=3$
and $\pt(T)=\{u_1,u_2,u_3\}$.
Define
$$k_j:=d(u_j,v)~~~j=1:3.$$
By Lemma \ref{equiv},
\[k_j\equiv q\,{\rm{mod}}\,(2q+1)~~~j=1:3.\]
Let $N_1$, $N_2$ and $N_3$ be such that
\[k_j=(2q+1)N_j+q~~~j=1:3.\]
Hence,
\begin{align} \label{kl+kt}
k_1+k_2+1=(2q+1)(N_1+N_2+1).
\end{align}
Let $b$ be an integer such that $b \in [0,q)$. Define
\begin{equation} \label{deltadefn}
\delta:=(2b+1)(N_1+N_2+1)~~\text{and}~~\gamma:=\frac{\delta}{k_1+k_2+1}. 
\end{equation}
Since $0 \leq b < q$,
\begin{equation} \label{deltae}
\delta < (2q+1) (N_1+N_2+1).
\end{equation}
By \eqref{kl+kt} and $\eqref{deltae}$,
$\delta$ is an integer in  the interval $(0,k_1 + k_2 +1)$.
Let the path $\p_{u_1u_2}$ in $T$ be
\[1:=u_1 \sim 2 \sim 3 \sim \cdots \sim k_1+1:=v \sim \cdots \sim k_1+k_2 \sim k_{1}+k_{2}+1:=u_2.\]
In view of Theorem \ref{th1}, 
\[\lambda:=2\left(1- \cos\left(\frac{{\delta \pi}}{{k_1+k_2+1}}\right)\right)\]
is an eigenvalue of $L_{\p_{u_1u_2}}$ and
there exists a vector $y=(y_1^{\delta},\dotsc,y_{k_1+k_2+1}^{\delta})'$ such that
$L_{\p_{u_1u_2}}y=\lambda y$, where 
\begin{equation}\label{eigvec}
 y_j^{\delta}:=\cos\left(\frac{\delta}{k_1+k_2+1} \left(j-\frac{1}{2} \right)\pi \right)~~~j=1:k_1+k_2+1.   
\end{equation}
 Since $k_1+1=v$, by \eqref{deltadefn} and \eqref{eigvec}, we obtain 
\[y_v^{\delta}=y_{k_1+1}^{\delta}=
        \cos \left(\gamma\left( k_1+\frac{1}{2} \right)\pi  \right).\]
As $k_1=(2q+1)N_1+q$ and $\gamma=\frac{2b+1}{2q+1}$,
       \[y_v^{\delta}= \cos \left( \frac{(2b+1)(2N_1+1)}{2}\pi \right).\]
        Thus,  $y_v^{\delta}=0$. Let $\p_{vu_3}$ be the path from $v$ to $u_3$.
        Then,
        $T=\p_{u_1u_2} \circ \p_{v u_3}$.
Define $f=(f_1,\dotsc,f_n)'$ by
\begin{equation}\label{xr}
f_r:=
\begin{cases}
y_r^{\delta} & r\in V(\p_{u_1u_2}) \\
0 & \mbox{else.}
\end{cases}    
\end{equation}
Since $y_v^\delta=0$, Lemma \ref{v1} implies that
$L_Tf=\lambda f$.
As $f_v=y_v^{\delta}$ and $y_v^{\delta}=0$, we get $f_v=0$.
Let $s \in V(T)$ be  such that $d(s,v) \equiv 0\,{\rm{mod}}\,(2q+1)$. 
We have two cases.
If $s\notin V(\p_{u_1u_2})$, then
by \eqref{xr}, $f_s=0$. If $s\in V(\p_{u_1u_2})$, then by Lemma \ref{lem1}, 
$y_s^{\delta}=0$, and again by
\eqref{xr},
\(f_s=0.\)
Thus, $f$ satisfies (a) and (b).


Applying the same argument to the path $\p_{u_2u_3}$,
we obtain a non-zero vector $g=(g_1,\dotsc,g_n)'$ such that
\begin{enumerate}
\item[(i)] $L_Tg=\lambda g$ 
\item[(ii)] $g_v=0$ and $g_i=0$ for all $i\notin V(\p_{u_2u_3})$ 
\item[(iii)] $g_s=0$ for all $s$ such that
$d(s,v)\equiv 0\,{\rm{mod}}\,(2q+1)$.
\end{enumerate}
Thus, $g$ satisfies (a) and (b).

We now claim that $f$ and $g$ are linearly independent.
Since
\[f_{u_1}=f_1=\cos\left(\frac{\gamma \pi}{2} \right) ~~\mbox{and} ~~0<\gamma<1,\]
we deduce that $f_1 \neq 0$.
Because
 $1=u_1 \notin V(\p_{u_2u_3})$, we see that
$g_1=0$.
Therefore, $f$ and $g$ are linearly independent.
Hence, the result is true for $p=3$.
We shall assume the result for trees with one major vertex and at most $p-1$ pendant vertices. 

Let $p>3$ and
$T$ be a tree with
$p$ pendant vertices. Let $u_1$ and $u_2$ be two distinct vertices in $\pt(T)$. Let the path $\p_{u_1u_2}$ be 
\[1:=u_1 \sim 2 \sim  \cdots \sim k_1+1:=v \sim \cdots \sim k_1+k_2+1:=u_2.\] 




Let $H$ be the component of $T\smallsetminus ({k_1})$ such that $v\in V(H)$. 
Then, $|\pt(H)|=p-1$. 
Since $p>3$, $H$ has a major vertex. 
By the induction hypothesis,
there exist linearly independent vectors $y^1,\dots,y^{p-2}$ such that
\begin{equation} \label{lh}
L_Hy^j=\lambda y^j,~~y^{j}_v=0~~\mbox{and}~~y^{j}_s=0~~~j=1:p-2
\end{equation}
for all $s \in V(H)$ such that
$d(s,v)\equiv 0\,{\rm{mod}}\,(2q+1)$.
Define $x^1,\dots,x^{p-2}$ in $\mathbb{R}^n$ as follows:
\begin{equation}\label{xi}
x^j_r:=
\begin{cases}
y^j_r & r\in V(H) \\
0 & \mbox{else.}
\end{cases}    
\end{equation}
As \[T=H\circ \p_{u_1v}~~\mbox{and}~~
y^j_v=0 ~~~ j=1:p-2;\] 
by Lemma \ref{v1}, 
\[L_T x^j=\lambda x^j~~~j=1:p-2.\]
In view of $\eqref{lh}$ and \eqref{xi},
\[x^j_v=0 ~~\mbox{and}~~x^j_s=0~~~j=1:p-2\]
for all $s\in V(T)$
such that $d(s,v) \equiv 0\,{\rm{mod}}\,(2q+1)$.
As $y^1,\dotsc,y^{p-2}$ are 
linearly independent, $x^1,\dotsc,x^{p-2}$
are linearly independent.

We now construct another eigenvector $x^{p-1}$. 
By Lemma \ref{equiv},
$$k_j\equiv q\,{\rm{mod}}\,(2q+1)~~j=1,2.$$ 
Hence, there exist integers $N_1$ and $N_2$ such that 
\[k_1=(2q+1)N_1+q~~\mbox{and}~~k_2=(2q+1)N_2+q.\] 
Therefore,
\begin{equation}\label{ksum}
    k_1+k_2+1=(2q+1)(N_1+N_2+1).
\end{equation}
Let $b$ be an integer in $[0,q)$. Put
\begin{equation} \label{deltadefn1}
\delta:=(2b+1)(N_1+N_2+1)~~\text{and}~~\gamma:=\frac{2b+1}{2q+1}. 
\end{equation}
By Theorem \ref{th1}, 
\[L_{\p_{u_1u_2}}(h_1^{\delta},\dotsc,h_{k_1+k_2+1}^{\delta})'=\lambda (h_1^{\delta},\dotsc,h_{k_1+k_2+1}^{\delta})',\]
where
\begin{equation}\label{hjdel}
 \lambda:=2\left(1-\cos (\gamma \pi) \right)~~\mbox{and}~~h_j^{\delta}:=\cos\left(\frac{\delta}{k_1+k_2+1} \left(j-\frac{1}{2} \right)\pi \right).   
\end{equation}
In particular, 
\[h_v^{\delta}=h_{k_1+1}^{\delta}=
        \cos \left(\gamma\left( k_1+\frac{1}{2} \right)\pi  \right).\]
Since $k_1=(2q+1)N_1+q$, we have
       $$h_v^{\delta}= \cos \left( \frac{(2b+1)(2N_1+1)}{2}\pi \right).$$
        Thus,  
        \begin{equation} \label{hv0}
            h_v^{\delta}=0.
        \end{equation}
        Define $x^{p-1}\in \rr^n$ by
\begin{equation}\label{p-1}
x^{p-1}_r:=
\begin{cases}
h_r^{\delta} & r\in V(\p_{u_1u_2}) \\
0 & \mbox{else.}
\end{cases}    
\end{equation}
Since $h_v^{\delta}=0$, using Lemma \ref{v1}, we obtain
$L_Tx^{p-1}=\lambda x^{p-1}$.
By \eqref{hv0} and \eqref{p-1},
$x^{p-1}_v=0$. Let $s\in V(T)$ be such that $d(s,v)\equiv 0\,{\rm{mod}}\,(2q+1)$. 
By Lemma  \ref{lem1}, 
$h_s^{\delta}=0$ for all $s\in V(\p_{u_1u_2})$.
Since $h_s^{\delta}=x_s^{p-1}$, we get $x_s^{p-1}=0$. If $s\notin V(\p_{u_1u_2})$, then \eqref{p-1} implies that $x_s^{p-1}=0$. Thus, $x^{p-1}$ satisfies condition (b).
By \eqref{xi},
\begin{equation}\label{xu1pi}
 x^j_{u_1}=0~~\mbox{for all}~~j=1:p-2.   
\end{equation}
By \eqref{hjdel} and \eqref{p-1},
\begin{equation*}
  x^{p-1}_{u_1}= \cos\left(\left( \frac{2b+1}{2q+1}\right)\frac{\pi}{2} \right). 
\end{equation*}
Since $0<\frac{2b+1}{2q+1}<1$, 
\begin{equation}\label{xjup-1}
  x^{p-1}_{u_1}\neq 0. 
\end{equation}
By \eqref{xu1pi} and \eqref{xjup-1}, $x^1,\dotsc,x^{p-1}$ are linearly independent.
The proof is complete.
\end{proof}

\end{document}
